% AA-46 | Picard--Vessiot i różniczkowa teoria Galois.
% Opracowanie ze źródła sources/pl/galois(2).tex i literatury:
% M. van der Put, M. Singer, Galois Theory of Linear Differential Equations.
% Twierdzenia głębokie (istnienie i jednoznaczność PV, odpowiedniość
% Galois, algebraiczność grupy, kryterium liouvillowskie) są JAWNIE cytowane.
\chapter{Rozszerzenia Picarda--Vessiota i różniczkowa teoria Galois}
\chaptermark{Teoria Picarda--Vessiota}
\label{aa:ch:picard-vessiot}

Klasyczna teoria Galois przyporządkowuje równaniom algebraicznym
grupy automorfizmów ciał rozkładu. Dla liniowych równań
różniczkowych analogiczną rolę pełnią rozszerzenia
Picarda--Vessiota i ich grupy automorfizmów różniczkowych.
Derywacja ogranicza automorfizmy: obraz rozwiązania musi
pozostać rozwiązaniem \emph{tego samego} równania. Zobaczymy,
jak wyprowadza to reprezentację grupy w $\mathrm{GL}_n(C)$.

Odwołujemy się do konstrukcji i kryteriów z rozdziałów
\ref{aa:ch:differential-fields}--\ref{aa:ch:nonelementary-integrals}.
W przeciwieństwie do pełnego dowodu twierdzenia Liouville'a,
kilka fundamentalnych twierdzeń strukturalnych teorii
Picarda--Vessiota przytoczymy jawnie jako wyniki klasyczne.
Ich pełne dowody wymagają rozwiniętej algebry różniczkowej
i teorii grup algebraicznych. Dowody elementarnych
lematów i wszystkich bezpośrednich rachunków wykonamy tutaj.

\section{Układy liniowe i macierze fundamentalne}

\begin{definicja}[Układ liniowy i macierz fundamentalna]
\label{aa:def:pv-fundamental}
Niech $(K,D)$ będzie ciałem różniczkowym
charakterystyki $0$, $C=C_K$, $n\ge1$
i $A\in M_n(K)$. Przez układ
\begin{equation}\label{aa:eq:pv-system}
 DY=AY
\end{equation}
rozumiemy równanie dla wektorów kolumnowych
$y\in L^n$, gdzie $L/K$ jest rozszerzeniem
różniczkowym, a $Dy$ oznacza różniczkowanie
wszystkich współrzędnych. Macierz
$U\in\mathrm{GL}_n(L)$ spełniającą
$DU=AU$ nazywamy \emph{macierzą fundamentalną}.
Jej kolumny są liniowo niezależnymi
rozwiązaniami nad ciałem stałych $C_L$.
\end{definicja}

\begin{lemat}[Przestrzeń rozwiązań]
\label{aa:lem:pv-solutions-constants}
Niech $L/K$ będzie rozszerzeniem różniczkowym,
$A\in M_n(K)$ i $U\in\mathrm{GL}_n(L)$
spełnia $DU=AU$. Wówczas
\[
 \{y\in L^n:Dy=Ay\}=U(C_L)^n.
\]
W szczególności przestrzeń rozwiązań
jest $n$-wymiarowa nad $C_L$.
Jeśli $V\in\mathrm{GL}_n(L)$ jest inną
macierzą fundamentalną, to $V=UM$
dla dokładnie jednej macierzy
$M\in\mathrm{GL}_n(C_L)$.
\end{lemat}
\begin{proof}
Z reguły różniczkowania odwrotności macierzy
$D(U^{-1})=-U^{-1}(DU)U^{-1}=-U^{-1}A$.
Dla rozwiązania $y$ otrzymujemy
\[
 D(U^{-1}y)=-U^{-1}Ay+U^{-1}Dy=0,
\]
więc współrzędne $U^{-1}y$ należą do $C_L$.
Odwrotnie, jeśli $c\in(C_L)^n$,
to $D(Uc)=(DU)c+U(Dc)=AUc$.
Stąd pierwsza równość.
Stosując ją do kolumn $V$,
otrzymujemy $V=UM$ z $M\in M_n(C_L)$,
a odwracalność $U,V$ wymusza
$M\in\mathrm{GL}_n(C_L)$. Jednoznaczność
wynika z odwracalności $U$.
\end{proof}

\begin{stwierdzenie}[Równanie Abela dla wyznacznika]
\label{aa:st:pv-abel-determinant}
Niech $U\in\mathrm{GL}_n(L)$ spełnia
$DU=AU$ dla $A\in M_n(K)$. Wówczas
\[
 D(\det U)=\operatorname{tr}(A)\det U.
\]
W szczególności, jeśli $\operatorname{tr}(A)=0$,
to $\det U\in C_L^\times$.
\end{stwierdzenie}
\begin{proof}
Wzór na pochodną wyznacznika
można otrzymać przez rozwinięcie
multiliniowe po kolumnach,
a po pomnożeniu przez $U^{-1}$
przyjmuje on postać
$D(\det U)=(\det U)\operatorname{tr}(U^{-1}DU)$.
Ponieważ $DU=AU$ i
$\operatorname{tr}(U^{-1}AU)=\operatorname{tr}(A)$,
otrzymujemy tezę. Gdy ślad jest zerowy,
pochodna wyznacznika znika; wyznacznik
jest niezerowy, bo $U$ jest odwracalna.
\end{proof}

\begin{przyklad}[Równanie skalarne jako układ]
\label{aa:prz:pv-companion}
Niech $a_0,\ldots,a_{n-1}\in K$.
Równanie
$D^ny+a_{n-1}D^{n-1}y+\cdots+a_0y=0$
jest równoważne układowi
$Dz=Az$ dla
$z=(y,Dy,\ldots,D^{n-1}y)^{\mathsf T}$,
gdzie $A$ ma jedynki na nadprzekątnej,
zera w pozostałych wierszach poza ostatnim,
a ostatni wiersz jest równy
$(-a_0,-a_1,\ldots,-a_{n-1})$.
Ostatnie równanie układu jest dokładnie
podanym równaniem skalarnym, a pozostałe
wiążą kolejne pochodne. Dzięki temu
wystarcza teoria macierzowa.
\end{przyklad}

\section{Różniczkowe pierścienie proste i rozszerzenia Picarda--Vessiota}

\begin{definicja}[Pierścień różniczkowy i ideał różniczkowy]
\label{aa:def:pv-differential-ideal}
Niech $R$ będzie przemiennym pierścieniem
z jedynką i derywacją $D:R\to R$.
Ideał $I\subseteq R$ nazywamy
\emph{różniczkowym}, jeśli $D(I)\subseteq I$.
Pierścień $R\ne0$ jest
\emph{różniczkowo prosty}, jeżeli
jego jedynymi ideałami różniczkowymi
są $(0)$ i $R$.
Prostota różniczkowa nie oznacza,
że $R$ musi być ciałem.
\end{definicja}

\begin{definicja}[Pierścień i rozszerzenie Picarda--Vessiota]
\label{aa:def:pv-extension}
Niech $(K,D)$ będzie ciałem różniczkowym
charakterystyki $0$ o algebraicznie
domkniętym ciele stałych $C=C_K$.
Dla $A\in M_n(K)$
\emph{pierścieniem Picarda--Vessiota}
dla układu~\eqref{aa:eq:pv-system}
nazywamy różniczkowo prostą
$K$-algebrę różniczkową $R$,
dla której istnieje
$U\in\mathrm{GL}_n(R)$ spełniająca
$DU=AU$ oraz
\[
 R=K[U_{ij},(\det U)^{-1}:
          1\le i,j\le n].
\]
Jest to zwykła algebra generowana
przez elementy macierzy i odwrotność
wyznacznika, a nie dowolna algebra
różniczkowa z dodatkowymi generatorami.

\emph{Ciałem Picarda--Vessiota} nazywamy
$E=\operatorname{Frac}(R)$, gdy $R$
jest dziedziną całkowitości.
Rozszerzenie $E/K$ nazywamy
\emph{rozszerzeniem Picarda--Vessiota}
dla rozważanego układu.
\end{definicja}

W definicji algebraiczna domkniętość $C$
jest założeniem, a nie automatyczną
własnością dowolnego ciała różniczkowego.
Pozwala sformułować standardową
teorię istnienia, jednoznaczności
i grup algebraicznych bez subtelności
związanych z rozszerzaniem stałych.

\begin{twierdzenie}[Istnienie i jednoznaczność -- wynik klasyczny]
\label{aa:tw:pv-existence-uniqueness}
Niech $(K,D)$ będzie ciałem różniczkowym
charakterystyki $0$ o algebraicznie
domkniętym ciele stałych $C_K=C$.
Dla każdego $n\ge1$ i każdego
$A\in M_n(K)$ istnieje pierścień
Picarda--Vessiota $R$.
Jest on dziedziną całkowitości,
a $E=\operatorname{Frac}(R)$ spełnia
$C_E=C$. Każde dwa takie rozszerzenia
są izomorficzne jako rozszerzenia
różniczkowe nad $K$
(nie kanonicznie).
\end{twierdzenie}

\begin{uwaga}[Co zawiera cytowany wynik?]
\label{aa:uw:pv-existence-cited}
To twierdzenie przytaczamy w klasycznej
postaci (van der Put--Singer,
\emph{Galois Theory of Linear Differential
Equations}, rozdział o teorii
Picarda--Vessiota).
Zarys konstrukcji jest następujący.
Na pierścieniu
$S=K[X_{ij},(\det X)^{-1}]$,
gdzie $X=(X_{ij})$ jest macierzą
formalnych zmiennych, rozszerzamy $D$
przepisem $DX=AX$.
Wybieramy maksymalny właściwy
ideał różniczkowy $I$.
Iloraz $S/I$ jest różniczkowo prosty
i zawiera fundamentalną macierz
obrazów $X_{ij}$.
Głębsze lemata algebry różniczkowej
gwarantują, że w charakterystyce
$0$ taki iloraz jest dziedziną
całkowitości i nie ma nowych stałych,
gdy $C$ jest algebraicznie domknięte.
Jednoznaczność wymaga ponadto
porównania dwóch pierścieni
fundamentalnych i ich
różniczkowych ideałów.
\emph{Właśnie te kroki są tu cytowane,
a nie przedstawiane jako dowiedzione.}
\end{uwaga}

\begin{uwaga}[Dlaczego nie wystarcza dowolne rozwiązanie?]
\label{aa:uw:pv-new-constants}
Dla równania $Dy=0$ i ciała
$K=\mathbb C(x)$ można dołączyć
przestępny element $t$ z $Dt=0$.
Nie jest to rozszerzenie
Picarda--Vessiota tego równania:
powstała nowa stała $t$, a
minimalne ciało rozwiązania to
już $K$. Sam fakt występowania
rozwiązań nie gwarantuje minimalności
ani zachowania ciała stałych.
\end{uwaga}

\section{Różniczkowa grupa Galois i jej reprezentacja}

\begin{definicja}[Różniczkowa grupa Galois]
\label{aa:def:pv-group}
Niech $E/K$ będzie rozszerzeniem
Picarda--Vessiota układu $DY=AY$.
\emph{Różniczkową grupą Galois}
nazywamy grupę
$G=\operatorname{Aut}_D(E/K)$
wszystkich automorfizmów ciała $E$,
które pozostawiają punktowo $K$
i komutują z derywacją:
$D(\sigma z)=\sigma(Dz)$
dla każdego $z\in E$.
\end{definicja}

\begin{twierdzenie}[Macierze automorfizmów]
\label{aa:tw:pv-automorphism-matrices}
Niech $(K,D)$ ma charakterystykę $0$
oraz algebraicznie domknięte ciało
stałych $C$, niech $E/K$ będzie
rozszerzeniem Picarda--Vessiota
dla $A\in M_n(K)$ i niech
$U\in\mathrm{GL}_n(E)$ będzie
macierzą fundamentalną
generującą $E$ nad $K$.
Dla każdego $\sigma\in G$
istnieje dokładnie jedna
$M_\sigma\in\mathrm{GL}_n(C)$ taka, że
\begin{equation}\label{aa:eq:pv-action}
 \sigma(U)=UM_\sigma.
\end{equation}
Odwzorowanie
$\rho_U:G\to\mathrm{GL}_n(C)$,
$\rho_U(\sigma)=M_\sigma$,
jest różnowartościowym
homomorfizmem grup.
Jeśli zmienimy bazę rozwiązań
$U$ na $UQ$ z
$Q\in\mathrm{GL}_n(C)$,
to macierze grupy zamieniają się
na $Q^{-1}M_\sigma Q$.
\end{twierdzenie}
\begin{proof}
Z $\sigma|_K=\mathrm{id}_K$ oraz
$D\sigma=\sigma D$ otrzymujemy
$D(\sigma U)=\sigma(DU)=\sigma(AU)=A\sigma U$.
Macierz $\sigma U$ jest odwracalna,
więc Lemat~\ref{aa:lem:pv-solutions-constants}
z $C_E=C$ daje
$\sigma U=UM_\sigma$ dla dokładnie
jednej macierzy $M_\sigma\in\mathrm{GL}_n(C)$.

Jeśli $\sigma,\tau\in G$, to
$(\sigma\circ\tau)U=\sigma(UM_\tau)
=(\sigma U)M_\tau=UM_\sigma M_\tau$,
ponieważ automorfizmy pozostawiają
$C\subseteq K$ punktowo.
Stąd $M_{\sigma\tau}=M_\sigma M_\tau$.
Jeśli $M_\sigma=I_n$, to $\sigma$
pozostawia wszystkie elementy $U$
nieruchome; ponieważ
$E=K(U_{ij}:1\le i,j\le n)$,
otrzymujemy $\sigma=\mathrm{id}_E$.
Dla zmiany $U'=UQ$ mamy
$\sigma(U')=UM_\sigma Q
=U'Q^{-1}M_\sigma Q$.
\end{proof}

\begin{twierdzenie}[Algebraiczność grupy -- wynik klasyczny]
\label{aa:tw:pv-algebraic-group}
Przy założeniach
Twierdzenia~\ref{aa:tw:pv-automorphism-matrices}
obraz $\rho_U(G)$ jest
\emph{domkniętą w topologii Zariskiego}
podgrupą $\mathrm{GL}_n(C)$,
czyli liniową grupą algebraiczną nad $C$.
Ponadto
$\operatorname{trdeg}_K E=\dim G$.
\end{twierdzenie}

\begin{uwaga}[Zakres twierdzenia]
\label{aa:uw:pv-algebraicity-cited}
Algebraiczność obrazu i formuła
dla stopnia przestępności są
wynikami teorii Picarda--Vessiota,
których tutaj nie dowodzimy.
Poprzednie twierdzenie daje
\emph{samodzielnie tylko} różnowartościowy
homomorfizm do $\mathrm{GL}_n(C)$.
Nie każda podgrupa $\mathrm{GL}_n(C)$
jest algebraiczna; domkniętości
Zariskiego nie wolno wyprowadzać
wyłącznie z działania na macierzy.
\end{uwaga}

\begin{twierdzenie}[Odpowiedniość Galois -- wynik klasyczny]
\label{aa:tw:pv-correspondence}
Niech $E/K$ będzie rozszerzeniem
Picarda--Vessiota w charakterystyce
$0$ z algebraicznie domkniętym
ciałem stałych $C$, a
$G=\operatorname{Aut}_D(E/K)$.
Wówczas $E^G=K$.
Odwzorowania
\[
 H\longmapsto E^H,\qquad
 F\longmapsto\operatorname{Aut}_D(E/F)
\]
ustalają odwracającą zawieranie
bijekcję między domkniętymi
podgrupami algebraicznymi $H\le G$
a pośrednimi ciałami różniczkowymi
$K\subseteq F\subseteq E$.
Pośrednie rozszerzenie $F/K$
jest rozszerzeniem Picarda--Vessiota
wtedy i tylko wtedy, gdy
odpowiadająca mu podgrupa
$H$ jest normalna w $G$;
w tym przypadku
$\operatorname{Aut}_D(F/K)\simeq G/H$.
\end{twierdzenie}

\begin{uwaga}[Dlaczego wymagamy domkniętości?]
\label{aa:uw:pv-correspondence-cited}
To pełne twierdzenie przytaczamy
za teorią klasyczną, a nie dowodzimy
z samego homomorfizmu~\eqref{aa:eq:pv-action}.
Istotne jest ograniczenie do
\emph{podgrup algebraicznych domkniętych
w topologii Zariskiego}. Dla
dowolnej, niedomkniętej podgrupy
odwzorowanie $H\mapsto E^H$
nie musi być różnowartościowe:
podgrupa i jej domknięcie Zariskiego
mogą mieć to samo ciało elementów stałych.
\end{uwaga}

\section{Przykłady: grupy multiplikatywna, addytywna i cykliczna}

\begin{przyklad}[Równanie $Dy=y$: grupa multiplikatywna]
\label{aa:prz:pv-multiplicative}
Niech $K=\mathbb C(x)$, $Dx=1$, a
$Dy=y$. Rozwiązaniem jest $t=e^x$.
Na mocy Lematu
\ref{aa:lem:exp-polynomial-constant-field}
$t$ jest przestępny nad $K$
i $C_{K(t)}=\mathbb C$.
Pierścień $R=K[t,t^{-1}]$
jest różniczkowo prosty.
Udowodnijmy ten ostatni fakt.
Niech $I\ne(0)$ będzie
ideałem różniczkowym $R$.
Wybierzmy niezerowy element
$p=\sum_{j=0}^m a_jt^j\in I$
o najmniejszej liczbie
niezerowych wyrazów; po
pomnożeniu przez odwrotność
monomianu $a_kt^k$ możemy
przyjąć, że $p$ ma
wyraz stały równy $1$.
Wtedy $Dp$ także należy do $I$,
ale wyraz stały jego pochodnej
jest zerowy; minimalność
liczby wyrazów daje $Dp=0$.
Każdy niezerowy współczynnik
$a_j$, $j\ne0$, musiałby
spełniać $Da_j/a_j=-j$.
W $\mathbb C(x)$ jest to
niemożliwe: pochodna logarytmiczna
funkcji wymiernej nie może być
niezerową stałą. Zatem $p=1$,
czyli $I=R$.

Jest to więc pierścień
Picarda--Vessiota, $E=K(t)$.
Każdy automorfizm różniczkowy
spełnia $\sigma(t)=ct$ dla
$c\in\mathbb C^\times$, a
każdy taki wybór daje
automorfizm $K(t)$ komutujący
z derywacją. Stąd
$G\simeq\mathbb G_m(\mathbb C)
=(\mathbb C^\times,\cdot)$.
\end{przyklad}

\begin{przyklad}[Równanie logarytmiczne: grupa addytywna]
\label{aa:prz:pv-additive}
Niech $K=\mathbb C(x)$ i
$t=\log x$ będzie przestępnym
generatorem o pochodnej $Dt=1/x$.
Brak nowych stałych wynika z
Lematu~\ref{aa:lem:primitive-extension-constants},
bo funkcja wymierna nie może
mieć pochodnej $1/x$.
Elementy $1,t$ są rozwiązaniami
równania $D^2y+x^{-1}Dy=0$.
Jeśli $z=(y,Dy)^{\mathsf T}$,
to jest to układ
\[
 Dz=\begin{pmatrix}0&1\\0&-1/x\end{pmatrix}z,
 \qquad
 U=\begin{pmatrix}1&t\\0&1/x\end{pmatrix}.
\]
Sprawdzamy wprost $DU=AU$ i
$\det U=1/x\in K^\times$.
Pierścień $R=K[t]$ jest
różniczkowo prosty:
dla niezerowego ideału
różniczkowego $I$ wybieramy
unormowany wielomian
$p=t^m+\sum_{j<m}a_jt^j\in I$
o najmniejszym stopniu.
Ponieważ $Dt=1/x$,
pochodna $Dp$ ma stopień
co najwyżej $m-1$, więc
z minimalności musi być zerowa.
Dla $m\ge1$ współczynnik
przy $t^{m-1}$ daje
$Da_{m-1}+m/x=0$,
co jest niemożliwe dla
funkcji wymiernej $a_{m-1}$,
gdyż pochodne wymierne
mają zerowe residua.
Zatem $m=0$ i $I=R$.

Automorfizmy mają postać
$t\mapsto t+c$ dla
$c\in\mathbb C$ i wszystkie
takie przesunięcia istnieją.
W bazie macierzy $U$ grupa
jest dana przez
\[
 M_c=\begin{pmatrix}1&c\\0&1\end{pmatrix},
 \qquad c\in\mathbb C,
\]
więc $G\simeq\mathbb G_a(\mathbb C)$.
\end{przyklad}

\begin{przyklad}[Rozszerzenie algebraiczne: grupa $\mu_m$]
\label{aa:prz:pv-cyclic}
Niech $m\ge2$, $K=\mathbb C(x)$
oraz $t^m=x$. Jedyna
przedłużona derywacja
spełnia $Dt=t/(mx)$.
Wielomian $T^m-x$ jest
nierozkładalny nad $\mathbb C(x)$
na mocy kryterium Eisensteina
w pierścieniu $\mathbb C[x]$
względem elementu pierwszego $x$:
współczynniki poza wiodącym
są podzielne przez $x$,
współczynnik stały $-x$
nie jest podzielny przez $x^2$,
a współczynnik wiodący $1$
nie jest podzielny przez $x$.
Samo wykluczenie $m$-tej potęgi
funkcji wymiernej nie wystarczałoby
do udowodnienia nierozkładalności
dla złożonego $m$.
Stąd $[K(t):K]=m$.
Ponieważ $C=\mathbb C$
jest algebraicznie domknięte,
$C_{K(t)}=\mathbb C$.
Pierścień $R=K[t,t^{-1}]=K(t)$
jest ciałem, zatem prostym
pierścieniem różniczkowym.
Rozwiązania równania
$Dy=(mx)^{-1}y$ generują $R$.
Każdy automorfizm wysyła
$t$ na $\zeta t$ z
$\zeta^m=1$, a każdy taki
wybór istotnie jest
automorfizmem różniczkowym.
Grupa wynosi
\[
 G=\mu_m=\{\zeta\in\mathbb C^\times:
                       \zeta^m=1\}.
\]
Ma ona $m$ elementów i wymiar
algebraiczny $0$, zgodnie z
$\operatorname{trdeg}_K K(t)=0$.
\end{przyklad}

\begin{przyklad}[Równanie harmoniczne]
\label{aa:prz:pv-harmonic}
Dla $K=\mathbb C(x)$ rozważmy
$D^2y+y=0$. Po podstawieniu
$t=e^{ix}$ dostajemy rozwiązania
$t$ i $t^{-1}$. Macierz
\[
 U=\begin{pmatrix}t&t^{-1}\\
                   it&-it^{-1}\end{pmatrix}
\]
spełnia $DU=AU$ dla
$A=\left(\begin{smallmatrix}0&1\\-1&0\end{smallmatrix}\right)$,
a $\det U=-2i\ne0$.
Mamy $E=K(t)$,
a grupa działa przez
$t\mapsto ct$, $c\in\mathbb C^\times$,
czyli
$\rho_U(G)=\{
\operatorname{diag}(c,c^{-1}):
c\in\mathbb C^\times\}$.
W szczególności wyznacznik
macierzy grupy wynosi $1$.
Nie wolno skalować rozwiązań
$t$ i $t^{-1}$ niezależnie,
ponieważ ich iloczyn wynosi $1$.
\end{przyklad}

\section{Rozszerzenia liouvillowskie i rozwiązania równań liniowych}

\begin{definicja}[Rozszerzenie liouvillowskie]
\label{aa:def:pv-liouvillian}
Niech $(K,D)$ ma charakterystykę $0$
i ciało stałych $C_K=C$.
Rozszerzenie różniczkowe $L/K$
bez nowych stałych nazywamy
\emph{liouvillowskim}, jeżeli istnieje
skończona wieża
$K=K_0\subseteq\cdots\subseteq K_r=L$,
w której $K_j=K_{j-1}(\theta_j)$,
a dla każdego $j$ zachodzi
co najmniej jeden warunek:
$\theta_j$ jest algebraiczny
nad $K_{j-1}$,
$D\theta_j\in K_{j-1}$ albo
$\theta_j\ne0$ i
$D\theta_j/\theta_j\in K_{j-1}$.
W odróżnieniu od rozszerzeń
elementarnych z
Definicji~\ref{aa:def:elementary-differential-extension}
krok $D\theta_j\in K_{j-1}$
pozwala dołączyć \emph{dowolną}
funkcję pierwotną, niekoniecznie
logarytm elementu niższego piętra.
\end{definicja}

\begin{twierdzenie}[Kryterium liouvillowskie -- wynik klasyczny]
\label{aa:tw:pv-liouvillian-criterion}
Niech $E/K$ będzie rozszerzeniem
Picarda--Vessiota układu $DY=AY$,
gdzie $\operatorname{char}K=0$,
a $C_K$ jest algebraicznie domknięte.
Niech $G=\operatorname{Aut}_D(E/K)$,
a $G^\circ$ oznacza składową spójną
jedności grupy algebraicznej $G$.
Wówczas $E/K$ jest rozszerzeniem
liouvillowskim wtedy i tylko wtedy,
gdy $G^\circ$ jest grupą rozwiązalną.
Dla równania drugiego rzędu
$D^2y=qy$, $q\in K$, istnienie
niezerowego rozwiązania w
pewnym rozszerzeniu liouvillowskim
jest równoważne istnieniu
elementu algebraicznego $r$
nad $K$ spełniającego
\begin{equation}\label{aa:eq:pv-riccati}
 Dr+r^2=q.
\end{equation}
\end{twierdzenie}

\begin{uwaga}[Zakres cytowanego kryterium]
\label{aa:uw:pv-lie-kolchin}
Kryterium to wykorzystuje
twierdzenie Liego--Kołczyna
o trójkątyzacji spójnych
rozwiązalnych grup liniowych
oraz odpowiedniość Galois
z Twierdzenia~\ref{aa:tw:pv-correspondence}.
Nie dowodzimy tu ani tego
twierdzenia, ani ogólnej
równoważności z liouvillowskością.
Prosty kierunek ostatniego
kryterium można jednak
sprawdzić bezpośrednio:
gdy $r$ jest algebraiczne
i spełnia~\eqref{aa:eq:pv-riccati},
dołączamy $r$, a następnie
niezerowy element $y$ z
$Dy/y=r$; otrzymujemy
$D^2y/y=Dr+r^2=q$.
Dla odwrotnego kierunku
potrzebna jest teoria
grup algebraicznych.
\end{uwaga}

\begin{przyklad}[Nieelementarna, ale liouvillowska]
\label{aa:prz:pv-erf-liouvillian}
Z rozdziału~\ref{aa:ch:nonelementary-integrals}
wiemy, że funkcja
$F(x)=\int_0^x e^{-u^2}\,\mathrm du$
nie ma elementarnej funkcji pierwotnej.
Jest jednak liouvillowska
w przyjętym znaczeniu:
najpierw dołączamy
$t=e^{-x^2}$ z
$Dt/t=-2x$, potem
element $F$ z $DF=t$.
Przestępność $t$ i brak nowych
stałych dla $K(t)$ wykazano
w Lemacie~\ref{aa:lem:exp-polynomial-constant-field}.
Jeśli $DF=t$ ma już rozwiązanie
w $K(t)$, nie trzeba nic
dołączać; w przeciwnym przypadku
przestępny krok z $DF=t$
również nie wprowadza nowych
stałych na mocy
Lematu~\ref{aa:lem:primitive-extension-constants}.
Ponadto $D^2F+2xDF=0$.
Jest to konkretny przykład
rozwiązania liniowego równania,
które nie jest elementarne,
choć należy do wieży liouvillowskiej.
\end{przyklad}

\section{Równanie Airy'ego i grupa specjalna liniowa}

\begin{lemat}[Przeszkoda Riccatiego dla równania Airy'ego]
\label{aa:lem:pv-airy-riccati}
Równanie $Dr+r^2=x$,
gdzie $D=\mathrm d/\mathrm dx$,
nie ma rozwiązań algebraicznych
nad $\mathbb C(x)$.
\end{lemat}
\begin{proof}
Przypuśćmy, że istnieje
element algebraiczny $r$
spełniający równanie.
Lokalnie przy każdym skończonym
$x=a$ gałęzie funkcji
algebraicznej mają rozwinięcia
Puiseux w potęgach
$(x-a)^{1/e}$.

Najpierw sprawdźmy, że żadna
taka gałąź nie może być
rozgałęziona w skończonym punkcie.
Jeżeli $r$ jest tam ograniczona,
równanie $r'=x-r^2$
ma holomorficzną prawą stronę
i lokalne rozwiązanie
o zadanej wartości $r(a)$
jest holomorficzne
i jednoznaczne.
Również porównanie najniższej
niecałkowitej potęgi Puiseux
pokazuje, że taka potęga
nie może wystąpić:
różniczkowanie obniża ją
o $1$, podczas gdy
prawa strona jej nie wytwarza
na niższym poziomie.

Jeśli gałąź ma biegun,
kładziemy $u=1/r$.
Z równania wynika
$u'=1-xu^2$,
a $u(a)=0$.
Prawa strona jest holomorficzna,
więc poprzedni argument
zastosowany do $u$
wyklucza rozgałęzienie
również w tym przypadku.
Zatem wszystkie skończone
punkty są nierozgałęzione.

Otrzymujemy spójne skończone
nierozgałęzione przykrycie
płaszczyzny zespolonej,
odpowiadające wielomianowi
minimalnemu $r$.
Płaszczyzna zespolona jest
po prostu spójna, więc takie
przykrycie ma stopień $1$.
W tym kroku używamy klasycznego
twierdzenia o przykryciach
powierzchni Riemanna.
Wynika stąd $r\in\mathbb C(x)$.

Pokażmy, że nie istnieje
nawet rozwiązanie wymierne.
Jeżeli wielomianowa część
$r$ w nieskończoności ma
stopień $d\ge1$, to $r^2$
ma biegun rzędu $2d$,
podczas gdy $r'$ ma biegun
rzędu co najwyżej $d-1$,
więc $r'+r^2$ ma biegun
rzędu parzystego $2d\ge2$.
Funkcja $x$ ma tam
biegun rzędu $1$.
Jeśli $r$ jest regularna
w nieskończoności, to
$r'+r^2$ jest również
regularna, a $x$ nie.
W obu przypadkach sprzeczność.
\end{proof}

\begin{twierdzenie}[Równanie Airy'ego]
\label{aa:tw:pv-airy-group}
Niech $K=\mathbb C(x)$,
$D=\mathrm d/\mathrm dx$.
Równanie $D^2y-xy=0$
nie ma niezerowego rozwiązania
liouvillowskiego nad $K$,
a jego różniczkowa grupa
Galois jest izomorficzna
z $\mathrm{SL}_2(\mathbb C)$.
\end{twierdzenie}
\begin{proof}
Z Lematu~\ref{aa:lem:pv-airy-riccati}
oraz cytowanego kryterium
\ref{aa:tw:pv-liouvillian-criterion}
wynika pierwsze stwierdzenie.

Niech $E/K$ będzie rozszerzeniem
Picarda--Vessiota rozważanego
równania, a $U$ jego macierzą
fundamentalną. Dla towarzyszącego
układu $A=\left(\begin{smallmatrix}
0&1\\x&0\end{smallmatrix}\right)$
mamy $\operatorname{tr}(A)=0$.
Na mocy
Stwierdzenia~\ref{aa:st:pv-abel-determinant}
$\det U\in\mathbb C^\times$.
Wybierając bazę rozwiązań
o wyznaczniku $1$, dostajemy
$G\subseteq\mathrm{SL}_2(\mathbb C)$.

Gdyby składowa spójna
$G^\circ$ była rozwiązalna,
cytowane kryterium
\ref{aa:tw:pv-liouvillian-criterion}
zapewniałoby niezerowe
rozwiązanie liouvillowskie,
wbrew pierwszej części dowodu.
Zatem $G^\circ$ nie jest
rozwiązalna.
Każda właściwa spójna
algebraiczna podgrupa
$\mathrm{SL}_2(\mathbb C)$
jest jednak rozwiązalna:
jej algebra Liego jest
właściwą podalgebrą
trójwymiarowej algebry
$\mathfrak{sl}_2(\mathbb C)$,
ma więc wymiar co najwyżej
$2$, a każda algebra Liego
wymiaru co najwyżej $2$
jest rozwiązalna.
Stąd $G^\circ=\mathrm{SL}_2(\mathbb C)$,
a ponieważ $G$ leży
w $\mathrm{SL}_2(\mathbb C)$,
otrzymujemy
$G=\mathrm{SL}_2(\mathbb C)$.
\end{proof}

\begin{uwaga}[Granice samodzielności dowodu]
\label{aa:uw:pv-airy-deep-inputs}
Dowód lematu o Riccatim
wymaga standardowego faktu
o lokalnych rozwinięciach
Puiseux i prostego twierdzenia
o nierozgałęzionych przykryciach
$\mathbb C$. Przejście od
braku algebraicznych rozwiązań
Riccatiego do grupy Galois
korzysta jawnie z głębokiego
kryterium~\ref{aa:tw:pv-liouvillian-criterion}.
Nie należy przedstawiać
końcowego wyniku
$\mathrm{SL}_2(\mathbb C)$
jako konsekwencji samego
rachunku stopni funkcji
wymiernych.
\end{uwaga}

\section{Co pozostaje poza zakresem rozdziału?}

Odpowiedniość Galois wyjaśnia,
dlaczego struktura rozwiązań
układu liniowego jest związana
z grupą algebraiczną automorfizmów.
Nie daje ona jeszcze
algorytmu rozstrzygającego
dla każdego układu,
jaka jest jego grupa,
ani automatycznej procedury
znajdowania funkcji pierwotnych.

Algorytm Rischa dotyczy
elementarnego całkowania,
a algorytm Kovacica dotyczy
liouvillowskich rozwiązań
jednorodnych równań
drugiego rzędu o wymiernych
współczynnikach. Ich ogólne
schematy i dowody poprawności
wymagają osobnego rozdziału
i nie są objęte niniejszym
opracowaniem. Rozdział ten
kończy przygotowany wcześniej
blok AA-43--AA-46;
ewentualny AA-47 wymaga
osobnego zatwierdzenia zakresu.
